% Source: https://artofproblemsolving.com/community/c3h2139076_distinguishable_card_drawing
% Stable identity: aops-topic-2139076-q01
Write $I$ for an indistinguishable card. The first two cards decide which case we are in (and how long the sequence is), so the cases give different draw orders and we can add them.

\textbf{Two distinguishable cards first.} The first two cards are an ordered pair of different $D$'s: $3\cdot2=6$ ways. The remaining cards are one $D$ and three $I$'s, and Bob draws $3$ of them: either $III$ ($1$ way) or the remaining $D$ in one of $3$ positions among two $I$'s ($3$ ways). Total $6\cdot4=24$.

\textbf{Two indistinguishable cards first.} The start $II$ can happen in $1$ way. The remaining cards are $D_1,D_2,D_3$ and one $I$, and Bob draws $3$ of them: all three $D$'s in some order ($3!=6$ ways), or two $D$'s and the $I$ ($3\cdot2=6$ ordered choices of the $D$'s times $3$ positions for the $I$: $18$ ways). Total $1\cdot24=24$.

\textbf{One of each first.} The start is $D_iI$ or $ID_i$: $2\cdot3=6$ ways. The remaining cards are two $D$'s and two $I$'s, and Bob draws $2$ of them: two $D$'s in some order ($2$ ways), one $D$ and one $I$ ($2$ choices of $D$ times $2$ orders: $4$ ways), or $II$ ($1$ way). Total $6\cdot7=42$.

Altogether there are $24+24+42=\boxed{90}$ draw orders.
